但禁渔后的生态恢复并非总沿良性方向发展。部分水域的鱼类增长过快，水生植物和浮游动物受到过度摄食，水体随之变浑浊，底栖植被大面积减少，部分水域甚至出现大面积"水下荒漠"。与此同时，长江江豚已回升至约1249头，2024年重新发现长江鲟自然繁殖产卵场；而鳄雀鳝、巴西龟等外来物种已在鄱阳湖等流域支流存在稳定种群。十年禁渔确实恢复了鱼类资源，但恢复后的系统未必更加健康，营养级失衡、江湖连通不足和外来物种扩散依然存在。因此有必要用统一动力学模型把"资源恢复"与"系统健康"分开评估。
\cite{report2022,report2023}。
